\section{Orthogonal support without an extension hypothesis}
\label{sec:orthogonal-nonextendible}

\paragraph{Status.}
This is a separate new author derivation.  It does not change the
reviewed D-OL-06--08 text.  The complete written proof and invoked
source interfaces have passed the author check, and the exact
author supplement passed. Whole-proof reviews by /root and
/root/g34\_eightfolds also pass, with a separate exact reconstruction.
The numerical supplement checks constants and finite witnesses;
it is not a substitute for the universal arguments below.

Keep the notation $S=\Omega(V)$, $O=O(V)$, $d=\dim V$ and
$\nu(T)=a-\max_\lambda\dim\ker(T-\lambda I)$.
In characteristic two the natural quadratic space is nondefective
and $d$ is even.  All representations in this section are ordinary
complex representations of the natural group $S$.

\subsection*{D-OL-09. A bound for every ordinary character when $q\ge3$}

\begin{theorem}[D-OL-09]
\label{D-OL-09}
Let $d\ge22$, $q\ge3$, and let $\rho$ be any nontrivial
irreducible complex representation of $S$, of degree $a$.
Every nonscalar compatible normalizer $T$ of finite projective
order satisfies
\[
                          \nu(T)\ge7a/64.                 \tag{ON.1}
\]
No extension of $\rho$ to $O$ is assumed.

Consequently, in the main imaginary-quadratic six-dimensional
multiplicity-variation setting, a normal irreducibly acting central
product of selfdual ordinary orthogonal factors with $q_i\ge3$
and $d_i\ge22$ cannot give period-image dimension nine.  The
factors may have different fields and dimensions, and all compatible
normalizers, scalars and factor permutations are allowed.
\end{theorem}

\begin{proof}
The abelian radicals $U_e,U_f$ and their generation proof in
D-OL-06 use only a nontrivial irreducible $S$-module.  In particular
\[
       \langle U_e,U_f\rangle=S,\qquad
       b:=\dim\rho^{U_e}=\dim\rho^{U_f}\le a/2.             \tag{ON.2}
\]
The descriptions of their weights and the $\Omega(W)$-orbits
$\mathcal O_t$ also require no extension.  We first enlarge the
moving-rank range in that calculation.

Suppose $g\in S$ has prime order and natural moving rank
$2\le r\le4$.  A nondegenerate subspace containing
$\operatorname{im}(g-I)$ can be chosen of dimension at most $2r$,
by adjoining dual partners to the radical of that image.  Its
perpendicular space is fixed by $g$ and has dimension at least
$d-2r\ge14$.  It contains a hyperbolic pair.  In coordinates
adapted to such a pair, put $N=d-2\ge20$ and use the line radical
$U_e$ on $W$ from D-OL-06.  The moving rank on $W$ is still $r$.
The fixed space in $W$ has dimension $N-r$, and its polar form
has radical dimension at most $r$.  The same quadratic Gauss-sum
calculation gives
\[
 \frac{|\mathcal O_t\cap\ker(g-I)|}{|\mathcal O_t|}
 \le
 \frac{q^{N-r-1}+q^{\lceil N/2\rceil}}
      {q^{N-1}-q^{\lceil N/2\rceil}-1}
 \le F_N(q)<1/8.                                           \tag{ON.3}
\]
Indeed, for the fixed space the exponent in the error is at most
$(N-r)-(N-2r)/2=N/2$.  Thus it is independent of $r$ in this
range.  The final inequality is the universal $q\ge3$ estimate
(OS.13).  Prime-order weight orbits and (ON.2) now imply
\[
                         \nu(\rho(g))\ge7a/32.              \tag{ON.4}
\]

The full natural semilinear-normalizer description in D-OL-07
holds for all $q$, not only for the small fields in its final
application: its outer-order calculation (OS.16), field lifts,
centralizer proof and represented-central-quotient argument were
all stated and proved with arbitrary $q=p^f$.  The complete
published outer table used there is
\cite[Section~2.2.2, Table~2.1, p.~14]{HES:LiXia2022}.
Accordingly a compatible normalizer is implemented on $S$ by a
natural semilinear similitude $H$.

First suppose this implementing map is linear, $H=h$.
If $T$ is nonscalar, $h$ is not a natural scalar.  Some element
$u$ of $U_e\cup U_f$ then fails to commute with $h$, because these
two radicals generate $S$ and their natural centralizer is scalar.
Set
\[
             k=huh^{-1}u^{-1}\in S\setminus\{1\}.          \tag{ON.5}
\]
Both $u$ and $huh^{-1}$ have natural moving rank two, so
$\operatorname{rank}(k-I)\le4$.  Every power of $k$ fixes its
fixed space.  In particular a nonidentity prime-order power $g$
has moving rank at most four and cannot be a nonidentity central
scalar.  Its rank is at least two: a rank-one odd-characteristic
isometry is a reflection of determinant $-1$, whereas in even
characteristic rank one is excluded by the Dickson invariant.
Both statements are the exact rank-one facts used in D-OL-07.
Thus (ON.4) applies to $g$.

Here all group operators are in $S$, so the commutator identity
is exact, with no choice of an $O$-extension:
\[
 [T,\rho(u)]=\rho(k),\qquad
 \nu(\rho(k))\le\operatorname{rank}(\rho(k)-I)\le2\nu(T).
                                                               \tag{ON.6}
\]
The last inequality follows by subtracting a scalar from $T$ in
$T\rho(u)-\rho(u)T$.  Since taking powers cannot increase
projective support, (ON.4)--(ON.6) prove (ON.1) for every linear
normalizer.  The commutator need not have rank two; retaining
the rank-four possibility is essential.

For a general $T$, first take a nonscalar power of prime order
$l$ in the projective linear group.  Its support is no larger
than $\nu(T)$, so it suffices to treat that power, still denoted
$T$.  Its field component has order either one or $l$.
The former is the linear case just proved.  In the latter case
write $F=\mathbf F_q$, $F_0=\mathbf F_{q_0}$ and
\[
                        q=q_0^l.                           \tag{ON.7}
\]
Because $T^l$ is scalar, projective faithfulness, perfectness and
the semilinear-centralizer calculation give $H^l=cI$ for some
$c\in F_0^*$.  The norm rescaling and quadratic-form rescaling
of (OS.22)--(OS.23) apply without any change for this cyclic
extension of prime degree: after rescaling $H$ and $Q$,
\[
                  H^l=I,\qquad Q(Hv)=\sigma(Q(v)).           \tag{ON.8}
\]
The averaging and independence-of-field-automorphisms proof there
gives an $F$-basis of $H$-fixed vectors and a nondegenerate
quadratic form over $F_0$.  The dimension $d\ge22$ guarantees a
hyperbolic pair over $F_0$.  In the resulting line radical, $H$
acts on weight vectors by entrywise $\sigma$.

At most $q_0^N$ weight vectors are fixed.  In any nonzero orbit
the fixed proportion is therefore at most
\[
 \frac{q_0^N}{q^{N-1}-q^{\lceil N/2\rceil}-1}
 \le\frac{q_0^{\,l-N(l-1)}}{1-q^{-9}-q^{-19}}
 \le\frac{2^{-18}}{1-4^{-9}-4^{-19}}<1/4.                   \tag{ON.9}
\]
Here $N\ge20$, $l\ge2$, $q_0\ge2$ and $q\ge4$; in
particular $l-N(l-1)\le-18$.  Norm levels not fixed by $\sigma$
have no fixed vectors.  The weight-orbit inequality gives
$\nu(T)\ge(1/2)(1-1/l)(3/4)a\ge3a/16$, stronger than
(ON.1).  This proves the first assertion for every field degree.

For the source consequence, the tensor product of the selfdual
factor modules has a unique nondegenerate invariant bilinear
form up to scalar.  The inducing image preserves its line.
Its similitude character is nontrivial: otherwise its character
would be real-valued, while the induction formula puts the
original imaginary quadratic field inside its character field.
This is exactly the already proved field and form interface
in D-OL-05 and D-OL-08.

On tensor products, factor-preserving operators retain the
support fraction $7/64$ by convolution.  Nontrivial factor cycles
have support at least $1/4$, as proved in (OS.27).  At least two
inertias have nontrivial similitude multiplier, each with actual
moving rank at least $e/2$.  The actual genus-zero rank-nine
factorization has at least twelve genuine inertias and total
moving rank $2e+6$.  Therefore
\[
                      2e+6\ge e+10(7e/64),\qquad e\le64.
                                                               \tag{ON.10}
\]
The genus-zero assertion does not precede the degree estimate
without justification: $a_i\ge q_i^{d_i-4}(q_i-1)>64$ by
\cite[Theorem~3.7, pp.~165--167]{OL:GLT2024}, so the actual Euler
formula excludes genus at least two; genus one has at most six
genuine marks and image rank at most six.  Finally $e\ge a_i>64$
contradicts (ON.10).  No equality of the inducing and original
character fields, and no sufficiency of a source-dimension
inequality, has been used.
\end{proof}

\subsection*{D-OL-10. Nonextendible characters in even characteristic}

\begin{theorem}[D-OL-10]
\label{D-OL-10}
Let $q$ be even, $d\ge1136$, and let $\rho$ be a nontrivial
ordinary irreducible $S$-module of degree $a$ which does not extend
to $O$.  Every nonscalar compatible finite-projective-order
normalizer satisfies
\[
                           \nu(T)\ge11a/64.                 \tag{ON.11}
\]
An irreducibly acting normal central product of such factor
modules, including mixed even fields and compatible factor
permutations, cannot supply image rank nine for its actual
six-dimensional multiplicity variation.  This conclusion does
not require selfduality or a determinant assumption.
\end{theorem}

\begin{proof}
We first prove a support bound for every nonidentity element of
$S$, without an extension hypothesis.  The natural group $S$ has
trivial center in these dimensions in even characteristic and is
simple.  Thus a nontrivial irreducible representation is faithful.
Taking a nonidentity prime-order power reduces the claim to an
element $g$ of prime order $l$.

Suppose first that $g$ fixes a hyperbolic pair.  Let $r\ge2$ be
its moving rank on the complementary space $W$, of dimension
$N=d-2$.  The numerator estimate in (ON.3) remains valid for
arbitrary $r$, not only $r\le4$.  If $r\le N/2$, the restricted
polar rank is at least $N-2r$ and its Gauss error is at most
$q^{N/2}$.  If $r>N/2$, the trivial bound on that Gauss error
is $q^{N-r}\le q^{N/2}$.  In both cases the fixed proportion
is at most $F_N(q)$.  Uniformly for $N\ge20$ and $q\ge2$,
\[
 F_N(q)\le\frac{q^{-2}+q^{-9}}{1-q^{-9}-q^{-19}}
 \le\frac{2^{-2}+2^{-9}}{1-2^{-9}-2^{-19}}<5/16.             \tag{ON.12}
\]
Together with (ON.2) and prime-order weight orbits this yields
$\nu(\rho(g))\ge(1/2)(11/16)(a/2)=11a/64$.

Suppose that no hyperbolic pair is fixed.  If $l$ is odd, $g$
is semisimple and its fixed space is a nondegenerate quadratic
space.  A nondegenerate quadratic space of dimension at least
three over a finite field contains a hyperbolic pair.  Thus
$\dim\ker(g-I)\le2$.
If $l=2$, put $r=\operatorname{rank}(g-I)$.  The fixed space
has dimension $d-r$, and its polar radical has dimension at most
$r$, because its perpendicular space is $\operatorname{im}(g-I)$.
Its polar rank is therefore at least $d-2r$.  Polar rank at least
four would give a four-dimensional nondegenerate quadratic
subspace, hence a hyperbolic pair.  Consequently
\[
 \begin{array}{ll}
 d-\dim\ker(g-I)\ge d-2, & l\text{ odd},\\
 r\ge d/2-1,             & l=2.
 \end{array}                                                   \tag{ON.13}
\]
In either case the left side, which is the natural support
restricted to $\mu_2=\{1\}$ in characteristic two, is at least
$567$ when $d\ge1136$.

Apply the exact published estimate
\cite[Theorem~4.3.6, (4.3.2), pp.~1917--1918; complete proof
pp.~1918--1923; Remark~4.3.7, p.~1923]{HOS:LST2011} with
$N=567=7\cdot9^2$.  It is stated for rational points of the
simple simply connected classical group, including
$\operatorname{Spin}_d$.  Inflate $\rho$ along the standard
simply connected finite cover onto $S$.  The inflated character
is still nontrivial and irreducible, and its natural action is
the same orthogonal action used in (ON.13).  The applicable
cover is the ordinary nonexceptional one in these large
even-characteristic dimensions; see the cover discussion
\cite[p.~4; Chapter~5, Table~1, p.~35]{BI:HHM}, already used in
D-OL-01.  No faithfulness premise appears in the published
character estimate, so inflation causes no change in its scope.

For $2\le q<43$ the published $\gamma_q$ is $19/20$; for larger
$q$ it is $1/(\sqrt q-1)\le19/20$.
The support convention in its remark allows precisely $\mu_2$
for the spin group, so (ON.13) is the required support, not an
unjustified algebraic-closure support substitution.  We obtain
\[
 \frac{|\chi(g)|}{a}<(19/20)^9,\qquad
 \nu(\rho(g))>
 \frac{1-(19/20)^9}{2}a>11a/64.                              \tag{ON.14}
\]
The last step uses $|\chi(g)|\ge a-2\nu(\rho(g))$.
This proves (ON.11) for every nonidentity inner operator.

We now explain why nonextendibility is exactly the needed linear
normalizer condition.  In even characteristic $O/S$ is cyclic of
order two, with an involution representative given by a symmetry.
If $\rho$ is invariant under that outer coset, choose an intertwiner
$A$ for a symmetry $r$.  Schur's lemma gives $A^2=cI$.
Rescale $A$ so $A^2=I$; assigning $r\mapsto A$ then extends
$\rho$ to the semidirect product $O=S\rtimes\langle r\rangle$.
Conversely any extension makes its character invariant.  Hence a
nonextendible character has no compatible operator from the
nontrivial $O/S$ coset.

Every linear similitude is a scalar times an isometry in even
characteristic.  The complete natural normalizer description
therefore makes a compatible linear normalizer a scalar times
an inner operator.  It satisfies the just-proved bound.
For a general compatible $T$, take a projective-prime-order power
as in D-OL-09.  If its field part is nontrivial, the complete
norm/form descent and uniform weight count (ON.7)--(ON.9) give
$\nu(T)\ge3a/16>11a/64$.  This includes arbitrary field degrees
and mixed graph--field lifts; the lift is not presumed to be a
pure Frobenius map.  The scalar and inner cases are already
covered.  Thus (ON.11) holds for every compatible normalizer.

Convolution and the weighted factor-cycle bound give support at
least $11e/64$ for every nonscalar inertia on the tensor product.
A nonidentity scalar has actual moving rank $e$.  The factor
degree bound $a_i\ge q_i^{d_i-4}(q_i-1)>96$ excludes genus at
least two in the actual Euler formula.  Genus one again gives
at most six genuine marks and image rank at most six.
A rank-nine source would therefore have genus zero and at least
twelve genuine inertia points.  It would satisfy
\[
                         2e+6\ge12(11e/64),\qquad e\le96,
                                                               \tag{ON.15}
\]
contradicting $e\ge a_i>96$.  This source contradiction needs
no invariant bilinear form.
\end{proof}

\subsection*{D-OL-11. Strict remaining ordinary orthogonal bounds}

\begin{corollary}[D-OL-11]
\label{D-OL-11}
In the main imaginary-quadratic sixfold problem, assume the
inducing image has a transitive irreducibly acting normal central
product of ordinary quasisimple orthogonal factors.  Then:
\begin{enumerate}
\item if $q\ge3$ and the factor module is selfdual, its natural
dimension satisfies $d\le21$;
\item if $q$ is even, irrespective of selfduality, its natural
dimension satisfies $d\le1134$.
\end{enumerate}
These are necessary bounds and do not assert existence at any
remaining parameter.
\end{corollary}

\begin{proof}
The first assertion is D-OL-09.  For the second, suppose instead
that $d\ge1136$.  Transitivity preserves extension to $O$:
the natural semilinear similitudes implementing automorphisms
normalize $O$, so an $O$-extension pulls back to an $O$-extension.
Thus either all factor modules extend or none does.
In the latter case D-OL-10 gives the contradiction, without a
selfduality assumption.  In the former case every factor module
is selfdual, because all irreducible $O$-characters are real-valued,
as used at the exact primary locators in D-OL-04.
For $q\ge4$, D-OL-09 excludes this case already at $d\ge22$.
For $q=2$, the reviewed extendible-character theorem D-OL-08
provides the same exclusion, also at $d\ge22$.
Finally $d$ is even in characteristic two, so $d<1136$ means
$d\le1134$.
\end{proof}

\paragraph{Scope of the new inputs and certificate.}
The only new published numerical input beyond D-OL-06--08 is
the explicitly stated Larsen--Shalev--Tiep estimate and its
restricted-scalar support remark, at the full primary locators
given above.  The generation, Gauss-sum, normalizer, descent,
degree and actual Hurwitz interfaces are the already proved
ones, with their exact hypotheses retained.  The supplement
checks the rational constants and an actual rank-four commutator;
it does not infer a universal representation theorem from finite
matrix examples.
